\chapter{Registry and Transactional State Management}
\label{chap:registry-persistence}

The registry connects the GlobalPlatform management model to the persistent
state of oXiDe SE. Its architectural role is to determine which objects and
which versions of their state belong to the secure element at an observable
command boundary. Writing an object into persistent memory is not sufficient
to make that object part of this state: it must also be reachable through the
committed Registry index.

This chapter summarizes the relevant GlobalPlatform requirements and explains
how the oXiDe SE architecture can satisfy them. It defines the contract for
command execution, publication and recovery. Chapter~8 relates these
architectural mechanisms to the effective implementation.

%--------------------------------------------------
\section{GlobalPlatform Requirements Considered}
\label{sec:gp-registry-transaction-requirements}

OPEN maintains the GlobalPlatform Registry. Applications, including Security
Domains, and Executable Load Files require entries. The Registry holds
identity, association, privilege, lifecycle, resource and card-wide management
information. Its physical storage format is not prescribed
\gpref{gp-card}{§6.5}.

Completion and atomicity must be read at the boundary specified for each
operation:
\begin{itemize}
\item Loading, installation and Registry Update have process-specific
  completion rules. A sequence of commands is not, merely by being a
  sequence, one indivisible transaction \gpref{gp-card}{§9.3--§9.4}.
\item \texttt{PUT KEY} explicitly requires atomic handling of multiple keys
  or components in its data field, and of components transferred through a
  chain ending with the command that clears the continuation bit
  \gpref{gp-card}{§11.8.2.3.3}.
\item Deletion must not start if verification fails. Once started, it must
  complete in the current session or, after interruption, at least its
  Registry updates must complete in a subsequent session. The deletion
  response marks completion \gpref{gp-card}{§9.5--§9.5.4}.
\end{itemize}

These requirements distinguish atomic modification, process completion and
completion after interruption. The architecture must express all three,
without imposing one transaction boundary on every management command.

%--------------------------------------------------
\section{Logical Registry and Authority}
\label{sec:registry-logical-contract}

The oXiDe SE registry is one kernel-owned catalogue. Security Domains define
administrative authority and visibility over its objects; they do not maintain
independent authoritative registries. Its public management projection covers
packages, applications and Security Domains. Internally, the catalogue also
includes keys, data and serialized Rustlet state, so their changes can be
committed together. This internal extension does not turn every stored object
into a publicly enumerable GlobalPlatform Registry entry.

A logical entry identifies an object, its kind, its parent Security Domain
and its payload. The payload provides the kind-specific information:
\begin{itemize}
\item a package designates executable content;
\item an instance binds application state to a package;
\item a Security Domain combines its administrative state and authority
  attributes with its executable binding where applicable;
\item a key or data object holds content managed under its owning authority.
\end{itemize}

\subsection{Management Views}
\label{subsec:registry-get-store-data-fallback}

Management commands observe authorized projections of this catalogue.
\texttt{GET STATUS} describes managed Card Content, not private key values or
arbitrary application data. Data access remains subject to the owning
Security Domain's authority and interpretation: storing a generic data object
must not override administrative information derived from the platform or
Security Domain state.

The same visibility rules apply before and after a mutation. Transactional
publication changes the committed state; it does not grant a new way to bypass
authority checks or to discover tentative objects.

%--------------------------------------------------
\section{Registry Index and Versioned Persistent Objects}
\label{sec:persistence-memory-model}
\label{sec:append-only-transactional-registry}
\label{subsec:persistence-block-families}

The architecture separates the \emph{Registry index}, which identifies a
coherent set of entries and references, from their persistent payloads.
Executable content forms a Code entry; generic values and key material form
Data entries; serialized application or Security Domain state forms a
Registry-instance state entry. These names express logical roles, independent
of a particular binary encoding.

Let $R$ be a committed Registry index and let $G(R)$ denote the state obtained
by following its references. A modification prepares replacement payloads in
new storage and adjusts a working index in RAM. Existing payloads reachable
from $R$ remain intact. The working index may therefore refer to both
unchanged payloads and newly prepared versions.

The essential distinction is between \emph{physical persistence} and
\emph{logical commitment}. A replacement payload can already be written in
flash while remaining absent from $G(R)$. Neither ordinary object lookup nor
recovery may promote it to committed state merely because its bytes exist.
This separation permits incremental preparation without incremental exposure
of the operation's effects.

Publication establishes a replacement index $R'$ only after all the payloads
it requires are persistent and valid. Recovery must recognize a complete
publication, reject an incomplete one and preserve an unambiguous ordering of
committed versions. These are requirements on the persistent representation;
its encoding and validation mechanism belong to Chapter~8.

\subsection{Preservation and Reclamation}
\label{subsec:persistence-recycling-rule}

Until $R'$ is committed, $R$ and every payload needed to reconstruct $G(R)$
remain protected from reclamation. New payloads needed by the pending
operation must also remain protected. After commitment, payloads that no
longer belong to any required committed or pending state may be reclaimed.
The protection boundary must respect the storage medium's erase granularity.

This rule applies to an arbitrary set of changed entries. Updating an instance
and an associated data object, or deleting an entire authorized dependency
set, does not require separate publications. Unchanged payloads can be shared
between versions. Space exhaustion must prevent commitment of an incomplete
replacement rather than destroy the previous committed state.

%--------------------------------------------------
\section{One APDU as an Observable State Transition}
\label{sec:registry-architectural-requirements}
\label{subsec:persistence-commit-rule}

For the command/response execution model considered here, application activity
starts with receipt of a command and ends with transmission of its response.
There is no background application process advancing between two APDUs.
The external functional observer sees responses and the behavior of later
commands, including commands issued after a reset. Internal execution steps
need not each expose a separately consistent registry.

The architectural unit of commitment for an ordinary modifying APDU includes
\emph{all} its persistent effects: changes requested through kernel services,
changes to other managed objects and the final serialized state of the
Rustlet. A service invoked during execution participates in the current
command's working state; it does not publish an independent Registry index.

The command boundary is organized as follows:
\begin{enumerate}
\item The kernel resolves the authorized objects from the committed state and
  restores the Rustlet state needed for execution.
\item During execution, mutations affect the working registry and prepare new
  payload versions. Later accesses within the same command may use this
  working state, under the same authority checks.
\item After application processing returns, the kernel serializes the final
  Rustlet state and includes it with every other persistent effect selected
  by the command's outcome in the replacement index.
\item Once the complete replacement is ready, the kernel publishes it and
  makes the subsequent command's working view agree with that committed
  state. A response asserting successful completion follows commitment.
\end{enumerate}

Returning from application processing is therefore not itself the commit
point. Final serialization and publication are still part of the command's
completion. A read-only command need not publish an unchanged state. An error
response follows its own specified state transition: it may legitimately
retain an explicitly defined effect, but must not accidentally retain an
arbitrary prefix of unsuccessful processing. If a proposed mutation is
abandoned, its tentative RAM view must not be reused by a later command as
though it had committed.

\subsection{Why Internal Mutations Remain Unobservable}

Suppose a Rustlet changes a Data entry from $D_0$ to $D_1$ and its own
serialized state from $S_0$ to $S_1$. The old index designates $(D_0,S_0)$.
During execution, new storage may already contain $D_1$, while $S_1$ has not
yet been serialized. That physical situation does not change the committed
pair. The new index is published only when it can designate $(D_1,S_1)$.

\begin{center}
\begin{tabular}{|p{5.8cm}|p{7cm}|}
\hline
\textbf{Interruption point} & \textbf{State available after recovery} \\
\hline
During execution or final serialization &
The old pair $(D_0,S_0)$. Prepared payloads do not confer visibility. \\
\hline
During index publication &
Either $(D_0,S_0)$ or $(D_1,S_1)$, according to whether the replacement
publication is complete and valid. \\
\hline
After commitment, including during response transmission &
The new pair $(D_1,S_1)$, even if the host did not receive the full response. \\
\hline
\end{tabular}
\end{center}

The partial pair $(D_1,S_0)$ is not a recoverable outcome of this operation.
Thus a sequence of internal mutations can be understood as one abstract
transition at an APDU observation boundary. Loss of a response does not prove
absence of commitment; it creates uncertainty for the host about which
complete transition occurred, not permission to expose a partial transition.

This reasoning depends on all persistent effects using the common publication
boundary. Effects outside that boundary would need their own explicit
composition rule. It concerns functional observations through commands and
recovery; physical side-channel observations are outside this argument.

%--------------------------------------------------
\section{Multi-APDU Management Transactions}
\label{sec:registry-multi-apdu}

The oXiDe SE architecture explicitly handles two cases of multi-APDU
management transactions:
\begin{enumerate}
\item \textbf{Key loading through chained \texttt{PUT KEY} commands.}
  The kernel collects the key components covered by the operation and
  activates the complete result together. Multiple keys or components carried
  in a single command use the same atomic publication principle, but do not
  themselves constitute a multi-APDU sequence.
\item \textbf{Package loading through \texttt{INSTALL [for load]} followed by
  several \texttt{LOAD} commands.} The kernel prepares the executable content
  incrementally and publishes the package only when the final transfer and
  validation are complete.
\end{enumerate}

These are explicit management operations with command-specific control of
Registry index publication. Their architectural treatment does not require a
general transaction that a Rustlet can open in one APDU and commit in another.
GlobalPlatform Card Specification~2.3.1 defines no generic management APDUs
for bracketing arbitrary application commands with \texttt{BEGIN TRANSACTION}
and \texttt{COMMIT TRANSACTION} \gpref{gp-card}{§11}.

\subsection{Preparation and Atomic Activation}

Each APDU terminates with its own response. The kernel retains a
\emph{preparation state} for the next command; no process continues running
between commands. An intermediate response acknowledges progress, without
making the pending result operational. The architecture separates:
\begin{itemize}
\item committed operational objects, used by ordinary management and
  application requests;
\item pending payloads and the preparation context, reserved for authorized
  continuation of the management operation.
\end{itemize}

During key replacement, the old key remains the operational key at
intermediate response boundaries; a newly created key remains unavailable
until complete. The final command validates the full result, replaces the
operational references and discharges the preparation context in one
publication. No intermediate command independently activates a subset of the
key material covered by the operation.

The context identifies the operation, target, authority and expected progress.
Continuation commands are checked against it. Other commands, where permitted,
must neither use nor accidentally activate the tentative result. Final
publication must preserve any intervening committed changes or reject a
conflicting continuation, rather than publish a stale registry view.

\subsection{Package Loading}
\label{subsec:persistence-dynamic-package-loading}

\texttt{INSTALL [for load]} establishes the authorized preparation of a
package. Successive \texttt{LOAD} commands contribute portions of its Code
entry. Their responses confirm transfer progress; the committed operational
registry does not yet expose the package as available for installation or
execution. Persisting those portions does not give them execution authority.

The last \texttt{LOAD} completes the input. Once the whole executable content
has been validated, the kernel publishes the package reference in the Registry
index before acknowledging successful completion. An interruption before this
commitment leaves the package absent from the operational catalogue, even
though some or all of its bytes may already be persistent. Installation of an
instance and subsequent lifecycle changes have their own completion
boundaries; they are not implicitly part of this loading transaction.

\subsection{Persistence of Preparation State}

Where interruption permits abandonment, the kernel can retain preparation in
RAM and leave new payloads uncommitted. Reset then leaves the operational
result unchanged. Where continuation across reset is required, an internal
persistent preparation record can preserve progress without activating the
pending result. Its payloads remain protected from reclamation.

An intermediate index publication can therefore preserve preparation while
still designating the old operational objects. The key property is atomic
activation of the managed result, not the number of index versions. The
kernel determines the appropriate preparation and publication boundaries for
each of these two management sequences.

%--------------------------------------------------
\section{Transactional Contract for a Rustlet}
\label{sec:registry-rustlet-transactions}

A Rustlet receives atomic persistence at the APDU boundary. By completion of
the transmission of a response acknowledging a successful state-changing
command, its final serialized state and the associated managed-object changes
have been committed together. The kernel performs serialization and commitment
before that response; it does not wait until transmission has finished to
start persisting the result. If execution or transmission is interrupted,
recovery exposes a complete old or new state as described above.

\subsection{Joint Application and Security Domain Outcomes}
\label{subsec:registry-sd-application-outcomes}

A protected APDU involves both an application and its Security Domain (SD).
The SD processes the incoming command and protects the outgoing response;
the application executes between these stages. The final publication decision
is made after response production and protection, but before transmission.
An application returning normally is therefore necessary but not sufficient
for committing its changes.

The following matrix specifies the persistent versions selected at that
boundary, assuming successful serialization and persistence. ``Previous''
means the version committed before this APDU, not before the secure channel
was established. ``Updated SD'' includes only legitimate protocol evolution.

\begin{center}
\begin{tabular}{llll}
\hline
SD crash & Application crash & SD state & Application state \\
\hline
No  & No  & Updated  & Updated \\
No  & Yes & Updated  & Previous \\
Yes & No  & Previous & Previous \\
Yes & Yes & Previous & Previous \\
\hline
\end{tabular}
\end{center}

There is one Registry index publication, not separate commits for the SD and
application. In the second row, that index combines valid SD evolution with
the previous application state. The application's tentative managed-object
mutations are abandoned together with its serialized state. The healthy SD
may protect the application-error response and continue its secure channel.
Rolling back legitimate protocol counters or chaining state merely because
the application crashed would not implement this contract.

An SD crash instead invalidates the complete APDU transition: neither
participant's tentative persistent changes may be published. Its channel is
closed and its volatile session is discarded. In particular, a crash while
protecting a successfully produced application response must still abandon
the application's staged changes. Both participants can crash if the SD fails
while protecting the application's crash response. If the SD crashes during
incoming processing, the application need not have executed at all. No
plaintext or partially protected response payload is emitted after an SD
crash; only the failure status may be returned without protection.

Internal restoration points do not constitute publication boundaries. A
successful later stage cannot revoke a previously established global abort.
However, an application-local abort does not by itself invalidate legitimate
SD protocol progress. A normally returned error status is distinct from a
crash and follows the operation's specified state-transition policy.

An unchanged selected state requires no new index publication. Serialization
or persistence failure must not expose a partially committed participant
set. These are atomicity and outcome-selection requirements; the matrix does
not by itself assert the complete set of database ACID properties.

\subsection{No Independent Intra-APDU Persistence Transaction}

The Rustlet's state is not serialized during application processing. Its
intermediate mutations belong to the working state of the current command,
and service-requested object mutations share that same publication boundary.
An intra-APDU transaction therefore has no independent meaning as a
\emph{persistence transaction} in this architecture. There is no intermediate
Rustlet state to commit separately: the APDU is the intrinsic functional
observation boundary of oXiDe SE.

A Rustlet may still save and restore temporary values, validate speculative
computation or handle an error during execution. Such actions determine the
final state of the command; they do not create additional externally visible
persistence boundaries. This contract concerns atomic persistence, not an
obligation for every error response to restore the state at command entry.
The response and resulting state must agree with the application's specified
behavior.

\subsection{Application Protocols for Rollback Across APDUs}

A Rustlet can build a protocol spanning several APDUs on top of this guarantee.
It retains the information needed for later validation or rollback as part of
its own serialized state. For example, its state may contain an accepted
value, a candidate value and a protocol phase:
\begin{enumerate}
\item A preparation command preserves the accepted value and initializes the
  candidate and phase. This complete preparation state is committed when the
  command succeeds.
\item Later commands update the candidate while preserving the accepted
  value. Each response boundary commits a consistent protocol state.
\item An application-defined validation command promotes the candidate to the
  accepted value and clears the preparation state. An application-defined
  cancellation command instead discards the candidate and keeps the accepted
  value. Each choice is one atomic APDU transition.
\end{enumerate}

If the protocol changes working values before final validation, it must keep
sufficient undo information to restore them. Associated object changes must
follow the same application protocol and participate in the relevant APDU's
common publication. Merely retaining old physical storage does not provide a
Rustlet with a rollback facility.

This rollback is a new application state transition, not a request to undo
past Registry index publications. Already transmitted responses cannot be
withdrawn. If intermediate results must remain unobservable, the Rustlet must
expose only accepted values while preparation is pending. It must also define
continuation, cancellation and repeated-request behavior after a reset or a
lost response. Atomic persistence ensures that its saved protocol state is
coherent at each boundary; the Rustlet defines what that state means.

The model consequently requires neither a generic transaction held open
across application commands nor an intra-APDU persistence API. The kernel owns
the two management sequences above; the Rustlet owns any additional
application protocol built from atomic APDU transitions.

%--------------------------------------------------
\section{Interrupted Deletion and Idempotent Recovery}
\label{sec:registry-deletion-recovery}

Atomic publication ensures that a cumulative deletion does not expose an
arbitrary subset of its removals. Its working index removes the complete
authorized dependency set, then publishes the resulting state together.
However, old-or-new recovery alone does not establish eventual completion of
an already engaged deletion.

To meet that additional obligation, the architecture needs a recoverable
engagement decision. One sufficient extension is a persistent deletion record
that binds the authorized target set to a completion obligation. It is
established before destructive work that cannot be abandoned. Before this
engagement, failed validation leaves deletion unstarted; afterwards, recovery
must retain the obligation until the required Registry effects are complete.

Continuation runs during recovery or command handling, without assuming
background activity between APDUs. It must complete the required transition
before ordinary commands can observe an inconsistent intermediate state.
Publishing the resulting index and discharging the completion obligation
must form one recoverable decision. Reclamation of obsolete storage must
preserve any information still needed for that decision.

Recovery must be idempotent under repeated interruptions. Already completed
removals remain completed; revisiting them must not fail merely because the
object is now absent, repeat an unrelated effect or target a newly created
object with a reused identity. The completion obligation therefore concerns
the originally authorized objects, not a fresh interpretation of a deletion
request against whatever objects happen to exist after restart.

This requirement adds forward completion to atomic publication. It does not
change the common Registry index mechanism: that mechanism can commit both
an operation's recovery state and its final result at the appropriate
boundaries.

%--------------------------------------------------
\section{Discussion: From Publication to an Observable Contract}

The architectural argument is that one index selects a coherent object graph,
while replacement payloads remain tentative until the relevant publication.
For one APDU, that graph combines service-requested mutations with final
Rustlet serialization. For the two multi-APDU management sequences, the kernel controls preparation
and indivisible activation. A Rustlet can compose atomic APDU transitions
into its own validation or rollback protocol. For an engaged
deletion, a recoverable obligation additionally ensures forward completion.

A formal model should relate these internal transitions to traces of APDU
responses and post-recovery behavior. Its proof obligations are complete
mediation of persistent effects, preservation of the previous committed graph,
unambiguous publication, separation of preparation from operational state,
and idempotent discharge of recovery obligations. These conditions explain
why the chosen architecture supports the required transactional behavior;
they do not constitute a completed formal proof or an assessment of the
implementation. That assessment belongs to Chapter~8.
